\begin{enumerate}
\def\labelenumi{\arabic{enumi}.}
\setcounter{enumi}{2}
\item
  \textbf{Weights, Measures, and Performance Data}
\end{enumerate}

\subsection{Physical Characteristics}

\begin{center}
{\small
\setlength{\tabcolsep}{6pt}
\begin{tabular}{>{\raggedright}p{5.2cm}l}
\toprule
\textbf{Parameter} & \textbf{Value} \\
\midrule
Form factor & PocketQube 1P (50$\times$50$\times$50 mm) \\
Mass (estimated) & $\sim$200 g \\
Stowed frontal area & $\sim$25 cm$^2$ \\
Deployed frontal area (SRAB) & $\sim$147 cm$^2$ \\
SRAB wing area & $\sim$122 cm$^2$ (4 wings, Asa2 geometry) \\
\bottomrule
\end{tabular}
}
\end{center}

% TODO: fill measured mass after final integration; the 200 g figure
% is the current estimate used in the SRAB simulation.

\subsection{Recovery Performance}

The SRAB recovery performance was evaluated through two drop-test
campaigns and computational simulation (Section 4):

\begin{center}
{\small
\setlength{\tabcolsep}{6pt}
\begin{tabular}{>{\raggedright}p{5.2cm}l}
\toprule
\textbf{Parameter} & \textbf{Value} \\
\midrule
Terminal velocity (predicted, Test 2) & 10.05 m/s \\
Rotation rate (Test 2) & 439 rpm \\
Cone angle (Test 2) & 11.9$^{\circ}$ \\
SCSM regulatory range & 20--45 m/s \\
Margin vs. lower bound & 50\% below \\
Impact energy dissipated by autorotation & 99.5\% \\
REC 10.2.1 notification (non-parachute) & Sent \\
\bottomrule
\end{tabular}
}
\end{center}

% TODO: update terminal velocity to the final selected configuration
% (Asa3, 2 wings) once the quantitative drop test (Test 4, ~200 m) is
% performed. The computational model currently predicts ~20 m/s for the
% Asa3 2-wing configuration, within the SCSM 20-45 m/s range.

\subsection{Power and Autonomy}

\begin{center}
{\small
\setlength{\tabcolsep}{6pt}
\begin{tabular}{>{\raggedright}p{5.2cm}l}
\toprule
\textbf{Parameter} & \textbf{Value} \\
\midrule
Battery & Li-Ion 6000 mAh \\
Average current draw & $\sim$145 mA \\
Predicted autonomy & $>$40 h \\
SCSM requirement & $\geq$4 h \\
Margin & $>$10$\times$ \\
\bottomrule
\end{tabular}
}
\end{center}

\subsection{Flight Sensors}

\begin{center}
{\small
\setlength{\tabcolsep}{6pt}
\begin{tabular}{>{\raggedright}p{3.2cm}ll}
\toprule
\textbf{Sensor} & \textbf{Type} & \textbf{Measurement} \\
\midrule
ICM-20602 & 6-axis IMU & acceleration + angular rate \\
BME280 & Environmental & pressure, temperature, humidity \\
NEO-8M & GPS & position, velocity, time \\
RFM95W & LoRa radio & 915 MHz telemetry link \\
\bottomrule
\end{tabular}
}
\end{center}

% TODO: add measured mass, measured current draw and final autonomy
% bench test results (with 20% margin evidence) as they become
% available.
